% Pakete -- bewusst schlank, damit der Bericht auch auf einer
% BasicTeX-Installation baut (siehe README.md).
%
% Zwei Engines, ein Quelltext: pdfTeX bekommt inputenc, XeTeX/LuaTeX nicht.
% fontenc[T1] gilt für beide, sonst fehlen den Schriften die Umlaute.

\usepackage{iftex}
\usepackage[T1]{fontenc}
\ifPDFTeX
  \usepackage[utf8]{inputenc}
\fi
\usepackage{lmodern}
\usepackage{microtype}

% XeTeX mit T1-Schriften setzt das direkte Zeichen ß als "SS".
% ��ber newunicodechar wird es auf \ss{} umgebogen -- wirkt auch unter pdfTeX.
\usepackage{newunicodechar}
\newunicodechar{ß}{\ss{}}

% Ohne diese Zeile bleibt die Sprache Englisch: die Klassenoption ngerman
% reicht bei article nicht, babel wird dadurch nicht geladen.
\usepackage[ngerman]{babel}

\usepackage[a4paper,left=3cm,right=2.5cm,top=2.5cm,bottom=2.5cm]{geometry}

\usepackage{graphicx}
\usepackage{booktabs}
\usepackage{longtable}
\usepackage{array}
\usepackage{caption}
\usepackage{enumitem}
\usepackage{amsmath}
\usepackage{xcolor}
\usepackage{csquotes}

\usepackage[
  backend=bibtex,
  style=authoryear,
  sorting=nyt,
  maxcitenames=2,
  maxbibnames=99,
  giveninits=true,
  uniquename=init,
  doi=true,
  url=true
]{biblatex}
\AtBeginBibliography{%
  \setlength{\bibitemsep}{2pt}%
  \setlength{\bibparsep}{0pt}%
}

\usepackage[hidelinks]{hyperref}

% Lange URLs und DOIs muessen umbrechen duerfen. Als unteilbare Box lief eine
% Adresse im Literaturverzeichnis 56 pt ueber den Rand (Bau vom 24.11.2026).
% biblatex erlaubt Umbrueche an Ziffern, Gross- und Kleinbuchstaben, ohne ein
% Trennzeichen zu setzen.
\DeclareFieldFormat{url}{URL\addcolon\space\url{#1}}
\setcounter{biburlnumpenalty}{100}
\setcounter{biburlucpenalty}{100}
\setcounter{biburllcpenalty}{100}
